% GENERATED FILE. Edit canonical lessons, then run npm run generate:books.
% canonical-source-hash: fnv1a64:c62ee3eaa055d4e0
% canonical-lessons: GE-C204-leder, GE-C204-holz, GE-C204-metall, GE-C204-kunststoff, GE-C204-koerper

\chapter{Leather, Wood, Metal, Plastic, The body}
\label{ch:ge-leather-wood-metal-plastic-the-body}
\hlchaptermodality{204}

\begin{chapteropening}
\textbf{By the end of this chapter:} \emph{I can say leather, wood, metal, plastic and the body.}

Complete the last lesson of chapter 204: I can say leather, wood, metal, plastic and the body.
\end{chapteropening}

\section[das Leder]{das Leder — leather}

\label{lesson:GE-C204-leder}



\begin{quote}
Before the new one: say the German for wool, then the German for silk.
\end{quote}

\subsection*{You'll want to know: das Leder}
\textbf{das Leder} — \textquotedblleft{}leather\textquotedblright{}.

\begin{cousinweb}
\textbf{Leder} $\leftarrow$ Old High German \emph{ledar} $\leftarrow$ Proto-Germanic \emph{*leþrą} — inherited, the same word as English \emph{leather}. (EtymDB 2.1)
\end{cousinweb}

\begin{grammarlens}[title={What you've built}]
One more word for this chapter.
\end{grammarlens}

\subsection*{Your turn}
\begin{itemize}
\raggedright
  \item \emph{Say it:} \emph{das Leder}
  \item \emph{Say it:} \emph{das Leder}, once more
  \item \emph{Say it:} say \emph{das Leder}
  \item \emph{Recall:} say the German for wool, then the German for silk, then say \emph{das Leder} again
\end{itemize}

\begin{tcolorbox}[breakable,colback=teal!4,colframe=teal!35!black,title={Before you move on}]
What does das Leder mean? (\textquotedblleft{}Leather\textquotedblright{}.) Say it once more.
\end{tcolorbox}
\section[das Holz]{das Holz — wood}

\label{lesson:GE-C204-holz}



\begin{quote}
Before the new one: say the German for silk, then the German for leather.
\end{quote}

\subsection*{You'll want to know: das Holz}
\textbf{das Holz} — \textquotedblleft{}wood\textquotedblright{}.

\begin{cousinweb}
\textbf{Holz} $\leftarrow$ Middle High German \emph{holz} $\leftarrow$ Old High German \emph{holz}, \textquotedblleft{}wood\textquotedblright{} — inherited, the same word as English \emph{holt}, \textquotedblleft{}a small wood\textquotedblright{}. (EtymDB 2.1)
\end{cousinweb}

\begin{grammarlens}[title={What you've built}]
One more word for this chapter.
\end{grammarlens}

\subsection*{Your turn}
\begin{itemize}
\raggedright
  \item \emph{Say it:} \emph{das Holz}
  \item \emph{Say it:} \emph{das Holz}, once more
  \item \emph{Say it:} say \emph{das Holz}
  \item \emph{Recall:} say the German for silk, then the German for leather, then say \emph{das Holz} again
\end{itemize}

\begin{tcolorbox}[breakable,colback=teal!4,colframe=teal!35!black,title={Before you move on}]
What does das Holz mean? (\textquotedblleft{}Wood\textquotedblright{}.) Say it once more.
\end{tcolorbox}
\section[das Metall]{das Metall — metal}

\label{lesson:GE-C204-metall}



\begin{quote}
Before the new one: say the German for leather, then the German for wood.
\end{quote}

\subsection*{You'll want to know: das Metall}
\textbf{das Metall} — \textquotedblleft{}metal\textquotedblright{}.

\begin{cousinweb}
\textbf{Metall} $\leftarrow$ Latin \emph{metallum} $\leftarrow$ Greek \emph{métallon}, \textquotedblleft{}mine, metal\textquotedblright{} — borrowed, like English \emph{metal}. EtymDB 2.1 has the German word but not its source.
\end{cousinweb}

\begin{grammarlens}[title={What you've built}]
One more word for this chapter.
\end{grammarlens}

\subsection*{Your turn}
\begin{itemize}
\raggedright
  \item \emph{Say it:} \emph{das Metall}
  \item \emph{Say it:} \emph{das Metall}, once more
  \item \emph{Say it:} say \emph{das Metall}
  \item \emph{Recall:} say the German for leather, then the German for wood, then say \emph{das Metall} again
\end{itemize}

\begin{tcolorbox}[breakable,colback=teal!4,colframe=teal!35!black,title={Before you move on}]
What does das Metall mean? (\textquotedblleft{}Metal\textquotedblright{}.) Say it once more.
\end{tcolorbox}
\section[der Kunststoff]{der Kunststoff — plastic}

\label{lesson:GE-C204-kunststoff}



\begin{quote}
Before the new one: say the German for wood, then the German for metal.
\end{quote}

\subsection*{You'll want to know: der Kunststoff}
\textbf{der Kunststoff} — \textquotedblleft{}plastic\textquotedblright{}.

\begin{cousinweb}
\textbf{Kunststoff} is the word for \textquotedblleft{}art, skill\textquotedblright{} + \emph{Stoff}, \textquotedblleft{}material\textquotedblright{}: material made by skill, not found in nature. \emph{Kunst} is built on \emph{können}, \textquotedblleft{}to be able\textquotedblright{}, inherited from Proto-Germanic \emph{*kunnaną} (EtymDB 2.1).
\end{cousinweb}

\begin{grammarlens}[title={What you've built}]
One more word for this chapter.
\end{grammarlens}

\subsection*{Your turn}
\begin{itemize}
\raggedright
  \item \emph{Say it:} \emph{der Kunststoff}
  \item \emph{Say it:} \emph{der Kunststoff}, once more
  \item \emph{Say it:} say \emph{der Kunststoff}
  \item \emph{Recall:} say the German for wood, then the German for metal, then say \emph{der Kunststoff} again
\end{itemize}

\begin{tcolorbox}[breakable,colback=teal!4,colframe=teal!35!black,title={Before you move on}]
What does der Kunststoff mean? (\textquotedblleft{}Plastic\textquotedblright{}.) Say it once more.
\end{tcolorbox}
\section[der Körper]{der Körper — the body}

\label{lesson:GE-C204-koerper}



\begin{quote}
Before the new one: say the German for metal, then the German for plastic.
\end{quote}

\subsection*{You'll want to know: der Körper}
\textbf{der Körper} — \textquotedblleft{}the body\textquotedblright{}.

\begin{cousinweb}
\textbf{Körper} $\leftarrow$ Middle High German \emph{korper}, from Latin \emph{corpus}, \textquotedblleft{}body\textquotedblright{} (EtymDB 2.1 has the Latin word) — borrowed. English \emph{corpse} and \emph{corps} share it.
\end{cousinweb}

\begin{grammarlens}[title={What you've built}]
One more word for this chapter.
\end{grammarlens}

\subsection*{Your turn}
\begin{itemize}
\raggedright
  \item \emph{Say it:} \emph{der Körper}
  \item \emph{Say it:} \emph{der Körper}, once more
  \item \emph{Say it:} say \emph{der Körper}
  \item \emph{Recall:} say the German for metal, then the German for plastic, then say \emph{der Körper} again
\end{itemize}

\begin{tcolorbox}[breakable,colback=teal!4,colframe=teal!35!black,title={Before you move on}]
What does der Körper mean? (\textquotedblleft{}The body\textquotedblright{}.) Say it once more.
\end{tcolorbox}
